\subsection{TENSOR PRODUCT OF ALGEBRAS OVER A FIELD}

Let K be a commutative field and E, F two algebras over K whose respective unit elements \(e, e'\) are non-zero. Then the homomorphisms \(\eta_{\mathbf{E}}: \mathbf{K} \to \mathbf{E}\) and \(\eta_{\mathbf{F}}: \mathbf{K} \to \mathbf{F}\) (§ 1, no. 3) are injections which allow us to identify K with a subfield of E (resp. F). The canonical homomorphisms \(u: \mathbf{E} \to \mathbf{E} \otimes_{\mathbf{K}} \mathbf{F}\) and \(v: \mathbf{F} \to \mathbf{E} \otimes_{\mathbf{K}} \mathbf{F}\) , defined by \(u(x) = x \otimes e'\) and \(v(y) = e \otimes y\) are injective (II, § 7, no. 9, Proposition 19) and allow us to identify E and F with subalgebras of \(\mathbf{E} \otimes_{\mathbf{K}} \mathbf{F}\) , both having as unit element the unit element \(e \otimes e'\) of \(\mathbf{E} \otimes_{\mathbf{K}} \mathbf{F}\) . In \(\mathbf{E} \otimes_{\mathbf{K}} \mathbf{F}\) , \(\mathbf{E} \cap \mathbf{F} = \mathbf{K}\) (II, § 7, no. 9, Proposition 19).

If \(\mathbf{E}'\) and \(\mathbf{F}'\) are subalgebras of \(\mathbf{E}\) and \(\mathbf{F}\) respectively, the canonical homomorphism \(\mathbf{E}' \otimes_{\mathbb{K}} \mathbf{F}' \to \mathbf{E} \otimes_{\mathbb{K}} \mathbf{F}\) is injective and allows us to identify \(\mathbf{E}' \otimes_{\mathbb{K}} \mathbf{F}'\) with the subalgebra of \(\mathbf{E} \otimes_{\mathbb{K}} \mathbf{F}\) generated by \(\mathbf{E}' \cup \mathbf{F}'\) (II, § 7, no. 7, Proposition 14).

PROPOSITION 6. Let E, F be two non-zero algebras over a commutative field, K, C (resp. D) a subalgebra of E (resp. F) and C' (resp. D') the centralizer of C in E (resp. D in F). Then the centralizer of C ⊗\_K D in E ⊗\_K F is C' ⊗\_K D'.

It all reduces to verifying that an element \(z = \sum_{i} x_{i} \otimes y_{i}\) of the centralizer of \(C \otimes_{K} D\) ( \(x_{i} \in F, y_{i} \in F\) ) belongs to \(C' \otimes_{K} D'\) ; we know that

\[
\mathrm{C} ^ {\prime} \otimes_ {\mathbf {K}} \mathrm{D} ^ {\prime} = (\mathrm{C} ^ {\prime} \otimes_ {\mathbf {K}} \mathrm{F}) \cap (\mathrm{E} \otimes_ {\mathbf {K}} \mathrm{D} ^ {\prime})
\]

(II, § 7, no. 7, Corollary to Proposition 14). The \(y_{i}\) may be assumed to be linearly independent over \(\mathbf{K}\) ; for all \(x \in \mathbb{C}\) , of necessity \((x \otimes e')z = z(x \otimes e')\) , that is \(\sum_{i} (xx_{i} - x_{i}x) \otimes y_{i} = 0\) , whence \(xx_{i} = x_{i}x\) for all \(i\) (II, § 3, no. 7, Corollary 1 to Proposition 7); hence of necessity \(x_{i} \in \mathbb{C}'\) for all \(i\) and therefore \(z \in \mathbb{C}' \otimes_{\mathbb{K}} \mathbb{F}\) ; it can similarly be shown that \(z \in \mathbb{E} \otimes_{\mathbb{K}} \mathbb{D}'\) , whence the proposition.

COROLLARY. If \(\mathbf{Z}\) and \(\mathbf{Z}'\) are the respective centres of \(\mathbf{E}\) and \(\mathbf{F}\) , the centre of \(\mathbf{E} \otimes_{\mathbb{K}} \mathbf{F}\) is \(\mathbf{Z} \otimes_{\mathbb{K}} \mathbf{Z}'\) .

Let E and F be two subalgebras of an algebra G over a commutative field K; suppose that every element of E commutes with every element of F. Then the canonical injections \(i: E \to G\) , \(j: F \to G\) define a canonical homomorphism \(h = i \otimes j: E \otimes_{K} F \to G\) (no. 2, Proposition 5) such that

\[
(i \otimes j) (x \otimes y) = x y \quad \text { for } x \in \mathrm{E}, y \in \mathrm{F}.
\]

DEFINITION 4. Given an algebra G over a commutative field K, two subalgebras E, F of G are said to be linearly disjoint over K if they satisfy the following conditions:

\begin{enumerate}
\def\labelenumi{(\arabic{enumi})}
\item
  every element of \(\mathbf{E}\) commutes with every element of \(\mathbf{F}\) ;
\item
  the canonical homomorphism of \(\mathbf{E} \otimes_{\mathbf{K}} \mathbf{F}\) into \(\mathbf{G}\) is injective.
\end{enumerate}

PROPOSITION 7. Let G be an algebra over a commutative field K and E, F two subalgebras of G such that every element of E commutes with every element of F. For E and F to be linearly disjoint over K, it is necessary and sufficient that there exist a basis of E over K which is a free subset of G for the right F-module structure on G. When this is so:

\begin{enumerate}
\def\labelenumi{(\roman{enumi})}
\item
  the canonical homomorphism \(h: \mathbf{E} \otimes_{\mathbb{K}} \mathbf{F} \to \mathbf{G}\) is an isomorphism of \(\mathbf{E} \otimes_{\mathbb{K}} \mathbf{F}\) onto the subalgebra of \(\mathbf{G}\) generated by \(\mathbf{E} \cup \mathbf{F}\) ;
\item
  \(\mathbf{E} \cap \mathbf{F} = \mathbf{K}\) ;
\item
  every free subset of E (resp. F) over K is a free subset of G with its right or left F-module (resp. E-module) structure.
\end{enumerate}

The condition of the statement is obviously necessary, since every basis of E over K is a basis of \(E \otimes_{K} F\) with its right F-module structure (II, §3, no. 7, Corollary 1 to Proposition 7). To see that the condition is sufficient, note that the image H of \(E \otimes_{K} F\) under h is the set of sums \(\sum_{i} x_{i} y_{i} = \sum_{i} y_{i} x_{i}\) in G, with \(x_{i} \in E\) and \(y_{i} \in F\) ; if \((a_{\lambda})\) is a basis of E over K, H is therefore also the submodule of the (right or left) F-module G, generated by \((a_{\lambda})\) . The condition of the statement therefore means that there exists a basis \((a_{\lambda})\) of E which is also a basis of the F-module H; it follows that h is injective. Assertion (iii) follows from the fact that every free subset of E is contained in a basis of E (II, §7, no. 1, Theorem 2).

COROLLARY 1. For the canonical homomorphism of \(\mathbf{E} \otimes_{\mathbf{K}} \mathbf{F}\) into \(\mathbf{G}\) to be bijective, it is necessary and sufficient that there exist a basis of \(\mathbf{E}\) over \(\mathbf{K}\) which is a basis of the (right or left) F-module \(\mathbf{G}\) .

COROLLARY 2. Let E, F be two subalgebras of G, of finite rank over K and such that every element of E commutes with every element of F. For E and F to be linearly disjoint over K, it is necessary and sufficient that the subalgebra H of G generated by E ∪ F be such that

\[
[ \mathrm{H}: \mathrm{K} ] = [ \mathrm{E}: \mathrm{K} ]. [ \mathrm{F}: \mathrm{K} ].\tag{12}
\]

This says that the rank over K of the surjective canonical homomorphism \(h: \mathbf{E} \otimes_{\mathbf{K}} \mathbf{F} \to \mathbf{H}\) is equal to the rank of \(\mathbf{E} \otimes_{\mathbf{K}} \mathbf{F}\) over \(\mathbf{K}\) , which is equivalent to saying that this homomorphism is bijective (II, § 7, No.~4, Proposition 9).

